

\section{Regime-Local Reduced-Order Models}
\label{sec:regime_local_representation}

The Regime-Adaptive Localized Mixture-of-Experts Reduced-Order Model
(RAL-MoE-ROM) represents a parametrized flow by a collection of local
reduced dynamical systems rather than by a single global coordinate system.
The construction is motivated by the fact that the steady, Hopf-onset, and
periodic regimes occupy geometrically and dynamically distinct parts of the
solution manifold.  We first define the regime cover and leakage-free data
partition, then construct the pressure-gauged local POD spaces, their
encoders, decoders, and scalers, and finally project the physical operators.
Only after all constituent objects have been introduced do we collect them
into a complete local chart.  History features, phase conventions, learned
updates, and native time integrators belong to the specialist dynamics and
are therefore defined in Section~\ref{sec:hprs_moe_specialists}, rather than
being treated as part of the POD representation itself.


\subsection{Parameterized flow dynamics and regime decomposition}
\label{subsec:parametric_regime_decomposition}

Let $\Omega\subset\mathbb{R}^{2}$ denote the spatial domain and let
$\mu\in\mathcal{M}$ be the parameter.  The
high-fidelity state at time $t$ is written as
\begin{equation}
  \boldsymbol{x}(t;\mu)
  =
  \bigl(\boldsymbol{u}(t;\mu),p^{\circ}(t;\mu)\bigr),
  \label{eq:full_state_definition}
\end{equation}
where $\boldsymbol{u}$ is the velocity field and $p^{\circ}$ is a
gauge-corrected pressure field.  The parameter domain is covered by three
regime-local sets,
\begin{equation}
  \mathcal{M}
  \subseteq
  \mathcal{M}_{S}\cup\mathcal{M}_{H}\cup\mathcal{M}_{P},
  \qquad
  r\in\mathcal{R}:=\{S,H,P\},
  \label{eq:regime_cover}
\end{equation}
associated with steady, hopf-onset, and periodic dynamics.  The cover is not
required to be disjoint.  Its adjacent overlaps are precisely the locations
at which more than one local model may be considered:
\begin{equation}
  \mathcal{O}_{SH}=\mathcal{M}_{S}\cap\mathcal{M}_{H},
  \qquad
  \mathcal{O}_{HP}=\mathcal{M}_{H}\cap\mathcal{M}_{P}.
  \label{eq:adjacent_overlaps}
\end{equation}
The certified topology does not contain a direct $S$--$P$ edge.  Hence,
physical proximity of two parameters alone does not authorize a cross-chart
prediction; adjacency and dynamical admissibility are both required.


For regime $r$, let
\begin{equation}
  \mathfrak{D}_{r}^{q}
  =
  \left\{
    \left(
      \mu_i,\,
      \{t_{i,n},\boldsymbol{u}_{i,n},p_{i,n}\}_{n=0}^{N_i-1}
    \right)
  \right\}_{i\in\mathcal{I}_{r}^{q}},
  \qquad
  q\in\{\mathrm{tra},\mathrm{va},\mathrm{te}\},
  \label{eq:trajectory_populations}
\end{equation}
denote its training, validation, and test sets.  The associated
parameter sets obey
\begin{equation}
  \mathcal{M}_{r}^{\mathrm{tr}}
  \cap\mathcal{M}_{r}^{\mathrm{va}}
  =
  \mathcal{M}_{r}^{\mathrm{tr}}
  \cap\mathcal{M}_{r}^{\mathrm{te}}
  =
  \mathcal{M}_{r}^{\mathrm{va}}
  \cap\mathcal{M}_{r}^{\mathrm{te}}
  =\varnothing .
  \label{eq:complete_re_isolation}
\end{equation}
All snapshots  belonging to one parameter trajectory
remain in the same set.  Every local mean, POD basis, and coefficient
scaler introduced below is estimated exclusively from
$\mathfrak{D}_{r}^{\mathrm{tr}}$. This complete-parameter,
complete-trajectory partition prevents a future segment of an initialization
trajectory from leaking into another set.



\subsection{Area-weighted local POD construction}
\label{subsec:local_pod_and_gauge}

Let $N_u$ and $N_p$ denote the numbers of velocity and pressure degrees of
freedom, respectively, on the common finite-volume mesh.  Let
$\boldsymbol{w}=(w_1,\ldots,w_{N_p})^{\mathsf T}$ collect the cell areas,
with $w_j>0$.  Since the pressure associated with an incompressible flow is
determined only up to an additive spatial constant, a pressure gauge must be
fixed before computing the pressure mean and constructing the corresponding
POD space.  We select the weighted zero-mean representative through the
linear map
\begin{equation}
\Gamma_p(q) = q-\frac{\boldsymbol{w}^{\mathsf T}q}
{\boldsymbol{w}^{\mathsf T}\boldsymbol{1}}\boldsymbol{1},
\qquad
p^\circ:=\Gamma_p(p).
\label{eq:pressure_gauge}
\end{equation}
The gauge-fixed pressure belongs to the discrete zero-mean space
\begin{equation}
\mathcal{Q}_{0,h} = \left\{
q\in\mathbb{R}^{N_p}:
\boldsymbol{w}^{\mathsf T}q=0
\right\},
\label{eq:discrete_zero_mean_pressure_space}
\end{equation}
and therefore satisfies
$\boldsymbol{w}^{\mathsf T}p^\circ=0$.  The operator $\Gamma_p$ removes the
constant-pressure nullspace and is an idempotent projection onto
$\mathcal{Q}_{0,h}$.

The finite-volume quadrature induces the discrete physical inner products
\begin{equation}
\langle \boldsymbol{v},\boldsymbol{z}\rangle_{M_u} = \boldsymbol{v}^{\mathsf T}M_u\boldsymbol{z},
\qquad
\langle q,s\rangle_{M_p} = q^{\mathsf T}M_p s,
\label{eq:weighted_inner_products}
\end{equation}
with the associated norms
\begin{equation}
\|\boldsymbol{v}\|_{M_u} = \left( \boldsymbol{v}^{\mathsf T}M_u\boldsymbol{v} \right)^{1/2},
\qquad
\|q\|_{M_p} = \left( q^{\mathsf T}M_pq \right)^{1/2}.
\label{eq:weighted_discrete_norms}
\end{equation}
Here $M_u$ and $M_p$ are the velocity and pressure quadrature matrices,
respectively.  For the two-dimensional collocated finite-volume
discretization and component-block ordering adopted in this work, they are
given by
\begin{equation}
M_p = \operatorname{diag}(w_1,\ldots,w_{N_p}),
\qquad
M_u = \operatorname{diag}(M_p,M_p),
\label{eq:velocity_pressure_mass_matrices}
\end{equation}
so that the same cell-area quadrature is applied independently to the two
velocity components.

For each regime $r\in\mathcal{R}$, the local POD spaces are constructed
exclusively from the snapshots contained in the regime-specific training
population $\mathfrak{D}_{r}^{\mathrm{tr}}$.  Let
\begin{equation}
N_r^{\mathrm{tr}} = \sum_{i\in\mathcal{I}_{r}^{\mathrm{tr}}}N_i
\label{eq:number_regime_training_snapshots}
\end{equation}
denote the total number of training snapshots in regime $r$.  After
enumerating the snapshots from all training trajectories by a single index,
the corresponding ensemble is written as
\begin{equation}
\left\{
\left(
\boldsymbol{u}^{(r)}_s,
p^{(r),\circ}_s
\right)
\right\}_{s=1}^{N_r^{\mathrm{tr}}},
\qquad
p^{(r),\circ}_s = \Gamma_p\left(p^{(r)}_s\right).
\label{eq:regime_training_snapshot_ensemble}
\end{equation}
The regime-local velocity and pressure means are
\begin{equation}
\overline{\boldsymbol{u}}^{(r)} = \frac{1}{N_r^{\mathrm{tr}}}
\sum_{s=1}^{N_r^{\mathrm{tr}}}
\boldsymbol{u}^{(r)}_s,
\qquad
\overline{p}^{(r),\circ} = \frac{1}{N_r^{\mathrm{tr}}}
\sum_{s=1}^{N_r^{\mathrm{tr}}}
p^{(r),\circ}_s.
\label{eq:local_training_means}
\end{equation}
Since $\Gamma_p$ is linear and every gauge-fixed pressure snapshot belongs
to $\mathcal{Q}_{0,h}$, the local pressure mean also satisfies
\begin{equation}
\boldsymbol{w}^{\mathsf T} \overline{p}^{(r),\circ} = 0.
\label{eq:local_pressure_mean_gauge}
\end{equation}

The centered velocity and pressure snapshot matrices are defined as
\begin{align}
X_u^{(r)}
&=
\left[
\boldsymbol{u}^{(r)}_1-\overline{\boldsymbol{u}}^{(r)}
\ \cdots
\boldsymbol{u}^{(r)}_{N_r^{\mathrm{tr}}}
-\overline{\boldsymbol{u}}^{(r)}
\right],
\label{eq:centered_velocity_snapshots} \\
X_p^{(r)}
&=
\left[
p^{(r),\circ}_1-\overline{p}^{(r),\circ}
\ \cdots
p^{(r),\circ}_{N_r^{\mathrm{tr}}}
-\overline{p}^{(r),\circ}
\right].
\label{eq:centered_pressure_snapshots}
\end{align}

The regime-local velocity POD basis is defined as an
$M_u$-orthonormal basis that minimizes the total projection error of the
centered velocity snapshots:
\begin{equation}
\Phi_u^{(r)}
\in
\underset{
\substack{
\Psi\in\mathbb{R}^{N_u\times d_u}\\
\Psi^{\mathsf T}M_u\Psi=I_{d_u}
}
}{\operatorname{arg\,min}}
\bigl\|
X_u^{(r)}
-
\Psi\Psi^{\mathsf T}M_uX_u^{(r)}
\bigr\|_{M_u,F}^{2}.
\label{eq:velocity_pod_minimization}
\end{equation}
Similarly, the regime-local pressure POD basis solves
\begin{equation}
\Phi_p^{(r)}
\in
\underset{
\substack{
\Psi\in\mathbb{R}^{N_p\times d_p}\\
\Psi^{\mathsf T}M_p\Psi=I_{d_p}
}
}{\operatorname{arg\,min}}
\bigl\|
X_p^{(r)}
-
\Psi\Psi^{\mathsf T}M_pX_p^{(r)}
\bigr\|_{M_p,F}^{2},
\label{eq:pressure_pod_minimization}
\end{equation}
where the weighted Frobenius norm is defined by
\begin{equation}
\|Y\|_{M,F}^{2} = \operatorname{tr}\left(Y^{\mathsf T}MY\right).
\label{eq:weighted_frobenius_norm}
\end{equation}
Consequently, the local POD spaces are optimal, in the corresponding
finite-volume approximations of the physical $L^2$ norms, among all linear
subspaces of dimensions $d_u$ and $d_p$.

The above minimization problems are solved through singular-value
decompositions of the quadrature-weighted snapshot matrices:
\begin{align}
M_u^{1/2}X_u^{(r)}
&=
U_u^{(r)}\Sigma_u^{(r)}{V_u^{(r)}}^{\mathsf T},
&
\Phi_u^{(r)}
&=
M_u^{-1/2}U_u^{(r)}(:,1{:}d_u),
\label{eq:weighted_velocity_svd} \\
M_p^{1/2}X_p^{(r)}
&=
U_p^{(r)}\Sigma_p^{(r)}{V_p^{(r)}}^{\mathsf T},
&
\Phi_p^{(r)}
&=
M_p^{-1/2}U_p^{(r)}(:,1{:}d_p).
\label{eq:weighted_pressure_svd}
\end{align}
The resulting local POD bases satisfy
\begin{equation}
\Phi_u^{(r)}
\in
\mathbb{R}^{N_u\times d_u},
\qquad
\Phi_p^{(r)}
\in
\mathbb{R}^{N_p\times d_p},
\label{eq:local_pod_bases}
\end{equation}
and are orthonormal with respect to the discrete physical inner products:
\begin{equation}
{\Phi_u^{(r)}}^{\mathsf T}M_u\Phi_u^{(r)} = I_{d_u},
\qquad
{\Phi_p^{(r)}}^{\mathsf T}M_p\Phi_p^{(r)} = I_{d_p}.
\label{eq:weighted_orthonormality}
\end{equation}
Moreover, because every centered pressure snapshot belongs to the
zero-mean space $\mathcal{Q}_{0,h}$, the retained pressure modes satisfy
\begin{equation}
\boldsymbol{w}^{\mathsf T}\Phi_p^{(r)} = \boldsymbol{0}^{\mathsf T}.
\label{eq:pressure_modes_gauge_consistency}
\end{equation}

Let $\sigma_{u,j}^{(r)}$ and $\sigma_{p,j}^{(r)}$ denote the singular values
of the weighted velocity and pressure snapshot matrices, respectively.  Their
squares measure the contributions of the corresponding POD modes to the
total weighted fluctuation energy.  The cumulative energy fractions retained
by the local reduced spaces are
\begin{equation}
\eta_u^{(r)}(d_u) = 
\frac{
\displaystyle
\sum_{j=1}^{d_u}
\bigl(\sigma_{u,j}^{(r)}\bigr)^2
}{
\displaystyle
\sum_{j=1}^{\rho_u^{(r)}}
\bigl(\sigma_{u,j}^{(r)}\bigr)^2
},
\qquad
\eta_p^{(r)}(d_p) = 
\frac{
\displaystyle
\sum_{j=1}^{d_p}
\bigl(\sigma_{p,j}^{(r)}\bigr)^2
}{
\displaystyle
\sum_{j=1}^{\rho_p^{(r)}}
\bigl(\sigma_{p,j}^{(r)}\bigr)^2
},
\label{eq:local_pod_energy_fractions}
\end{equation}
where $\rho_u^{(r)}$ and $\rho_p^{(r)}$ are the respective ranks of the
weighted snapshot matrices.  The corresponding minimum squared projection
errors are
\begin{equation}
\bigl\|
X_u^{(r)}
-
\Phi_u^{(r)}{\Phi_u^{(r)}}^{\mathsf T}
M_uX_u^{(r)}
\bigr\|_{M_u,F}^{2}
= \sum_{j=d_u+1}^{\rho_u^{(r)}}
\bigl(\sigma_{u,j}^{(r)}\bigr)^2,
\label{eq:velocity_pod_truncation_error}
\end{equation}
and
\begin{equation}
\bigl\|
X_p^{(r)}
-
\Phi_p^{(r)}{\Phi_p^{(r)}}^{\mathsf T}
M_pX_p^{(r)}
\bigr\|_{M_p,F}^{2}
= \sum_{j=d_p+1}^{\rho_p^{(r)}}
\bigl(\sigma_{p,j}^{(r)}\bigr)^2.
\label{eq:pressure_pod_truncation_error}
\end{equation}
Thus, the singular values provide both an energy-based interpretation of the
retained modes and an exact characterization of the training-set POD
truncation error.  The local means, POD bases, and mode dimensions are
determined exclusively from the training population; validation and test
snapshots do not enter any stage of the POD construction.

For an arbitrary velocity field $\boldsymbol{u}$, its regime-local POD
coefficient vector is obtained by the $M_u$-orthogonal projection
\begin{equation}
\boldsymbol{a}^{(r)} = {\Phi_u^{(r)}}^{\mathsf T}M_u
\bigl( \boldsymbol{u} - \overline{\boldsymbol{u}}^{(r)} \bigr).
\label{eq:velocity_encode}
\end{equation}
For an arbitrary pressure field $p$, the pressure gauge is first fixed and
the local pressure coefficients are then defined by
\begin{equation}
\boldsymbol{b}^{(r)} = {\Phi_p^{(r)}}^{\mathsf T}M_p
\bigl( \Gamma_p(p) - \overline{p}^{(r),\circ} \bigr).
\label{eq:pressure_encode}
\end{equation}
The associated affine POD reconstruction maps are
\begin{align}
\mathcal{D}^{u}_{r}(\boldsymbol{a}^{(r)})
&= \overline{\boldsymbol{u}}^{(r)} + \Phi_u^{(r)}\boldsymbol{a}^{(r)},
\label{eq:velocity_decode} \\
\mathcal{D}^{p}_{r}(\boldsymbol{b}^{(r)})
&= \overline{p}^{(r),\circ} + \Phi_p^{(r)}\boldsymbol{b}^{(r)}.
\label{eq:pressure_decode}
\end{align}
The compositions of the coordinate and reconstruction maps define the
orthogonal projections onto the corresponding affine reduced spaces.  In
particular,
\begin{equation}
\mathcal{P}^{u}_{r}(\boldsymbol{u}) = \overline{\boldsymbol{u}}^{(r)}
+ \Phi_u^{(r)}{\Phi_u^{(r)}}^{\mathsf T}M_u
\bigl( \boldsymbol{u} - \overline{\boldsymbol{u}}^{(r)} \bigr),
\label{eq:velocity_affine_pod_projection}
\end{equation}
while the gauge-consistent pressure projection is
\begin{equation}
\mathcal{P}^{p}_{r}(p) = \overline{p}^{(r),\circ}
+ \Phi_p^{(r)}{\Phi_p^{(r)}}^{\mathsf T}M_p
\bigl( \Gamma_p(p) - \overline{p}^{(r),\circ} \bigr).
\label{eq:pressure_affine_pod_projection}
\end{equation}
These maps are referred to below as the regime-local POD coordinate and
reconstruction maps.  The terms encoder and decoder are used only when
describing their role in the implementation of the learned specialist
models.

As a numerical preprocessing step for the subsequent learned models, the
POD coefficients are standardized separately within each regime:
\begin{equation}
\widetilde{\boldsymbol{a}}^{(r)} = 
\bigl( \boldsymbol{a}^{(r)} - \boldsymbol{\mu}_{a}^{(r)} \bigr)
\oslash
\boldsymbol{\sigma}_{a}^{(r)},
\qquad
\widetilde{\boldsymbol{b}}^{(r)} = 
\bigl( \boldsymbol{b}^{(r)} - \boldsymbol{\mu}_{b}^{(r)} \bigr)
\oslash
\boldsymbol{\sigma}_{b}^{(r)},
\label{eq:local_scalers}
\end{equation}
where $\oslash$ denotes componentwise division, and
$\boldsymbol{\mu}_{a}^{(r)}$, $\boldsymbol{\sigma}_{a}^{(r)}$,
$\boldsymbol{\mu}_{b}^{(r)}$, and
$\boldsymbol{\sigma}_{b}^{(r)}$ are computed exclusively from the POD
coefficients of $\mathfrak{D}_{r}^{\mathrm{tr}}$.  Because the snapshots are
centered before projection, the coefficient means vanish in exact
arithmetic; they are retained in \eqref{eq:local_scalers} to describe the
implemented preprocessing uniformly.

The standardized coefficients are used only within the learned
regime-specific models.  The Galerkin and Pressure--Poisson operators,
inter-chart coordinate transformations, and physical-field reconstructions
are formulated in terms of the unstandardized POD coefficients
$\boldsymbol{a}^{(r)}$ and $\boldsymbol{b}^{(r)}$.










\subsection{Chart-local POD--Galerkin and Pressure--Poisson operators}
\label{subsec:local_galerkin_pressure_poisson}

The regime-local reduced operators are obtained by projecting the
semi-discrete incompressible flow equations onto the local POD spaces
introduced in Section~\ref{subsec:local_pod_and_gauge}.  After elimination of
prescribed boundary degrees of freedom and incorporation of the lifting
associated with the parameter-dependent boundary data, the full-order
velocity equation is written in mass-inverted form as
\begin{equation}
  \dot{\boldsymbol{u}}
  =
  \boldsymbol{f}_{h}(\mu)
  +\mathcal{L}_{h}(\mu)\boldsymbol{u}
  -\mathcal{N}_{h}
   \bigl(\boldsymbol{u},\boldsymbol{u};\mu\bigr)
  -\nabla_{h}(\mu)p^{\circ},
  \qquad
  \operatorname{div}_{h}(\mu)\boldsymbol{u}=0.
  \label{eq:semidiscrete_incompressible_system}
\end{equation}
Here $\mathcal{L}_{h}(\mu)$ denotes the discrete linear momentum operator,
including the viscous contribution, and
$\mathcal{N}_{h}(\boldsymbol{v},\boldsymbol{z};\mu)$ denotes the bilinear
discrete convection operator associated with
$(\boldsymbol{v}\cdot\nabla)\boldsymbol{z}$.  The operators
$\nabla_h(\mu)$ and $\operatorname{div}_h(\mu)$ are the discrete gradient
and divergence, respectively, while $\boldsymbol{f}_{h}(\mu)$ collects the
body-force, boundary, and lifting contributions that are independent of the
instantaneous reduced coordinates.  The explicit dependence on $\mu$
allows for parameter-dependent boundary elimination, lifting functions, or
offline interpolation of full-order operators.  Equation
\eqref{eq:semidiscrete_incompressible_system} also fixes the sign convention
used throughout this subsection.

Using the affine POD reconstructions
\eqref{eq:velocity_decode}--\eqref{eq:pressure_decode}, the local
approximations in regime $r$ are written componentwise as
\begin{equation}
  \boldsymbol{u}^{(r)}_{\mathrm{ROM}}
  =
  \overline{\boldsymbol{u}}^{(r)}
  +\sum_{j=1}^{d_u}
   a_j^{(r)}\boldsymbol{\phi}^{u,(r)}_j,
  \qquad
  p^{\circ,(r)}_{\mathrm{ROM}}
  =
  \overline{p}^{(r),\circ}
  +\sum_{\ell=1}^{d_p}
   b_\ell^{(r)}\phi^{p,(r)}_\ell.
  \label{eq:local_pod_expansions_componentwise}
\end{equation}
Substitution of \eqref{eq:local_pod_expansions_componentwise} into the
momentum equation in \eqref{eq:semidiscrete_incompressible_system}, followed
by Galerkin projection with respect to the $M_u$ inner product, yields
\begin{equation}
  M^{u,(r)}\dot{\boldsymbol{a}}^{(r)}
  =
  \boldsymbol{c}_r(\mu)
  +A_r(\mu)\boldsymbol{a}^{(r)}
  +H_r(\mu):
   \bigl(
     \boldsymbol{a}^{(r)}
     \otimes
     \boldsymbol{a}^{(r)}
   \bigr)
  +C_r^p(\mu)\boldsymbol{b}^{(r)}.
  \label{eq:projected_local_momentum_general_mass}
\end{equation}
The reduced velocity mass matrix is
\begin{equation}
  \bigl[M^{u,(r)}\bigr]_{ij}
  =
  \left\langle
    \boldsymbol{\phi}^{u,(r)}_i,\,
    \boldsymbol{\phi}^{u,(r)}_j
  \right\rangle_{M_u}.
  \label{eq:local_reduced_velocity_mass}
\end{equation}
Because the local velocity POD basis is $M_u$-orthonormal,
$M^{u,(r)}=I_{d_u}$.  The remaining reduced operators are defined by
\begin{align}
  [\boldsymbol{c}_r(\mu)]_i
  &=
  \left\langle
    \boldsymbol{\phi}^{u,(r)}_i,\,
    \boldsymbol{f}_{h}(\mu)
    +\mathcal{L}_{h}(\mu)\overline{\boldsymbol{u}}^{(r)}
    -\mathcal{N}_{h}
     \bigl(
       \overline{\boldsymbol{u}}^{(r)},
       \overline{\boldsymbol{u}}^{(r)};\mu
     \bigr)
    -\nabla_{h}(\mu)\overline{p}^{(r),\circ}
  \right\rangle_{M_u},
  \label{eq:galerkin_constant_definition}\\
  [A_r(\mu)]_{ij}
  &=
  \left\langle
    \boldsymbol{\phi}^{u,(r)}_i,\,
    \mathcal{L}_{h}(\mu)\boldsymbol{\phi}^{u,(r)}_j
    -\mathcal{N}_{h}
     \bigl(
       \overline{\boldsymbol{u}}^{(r)},
       \boldsymbol{\phi}^{u,(r)}_j;\mu
     \bigr)
    -\mathcal{N}_{h}
     \bigl(
       \boldsymbol{\phi}^{u,(r)}_j,
       \overline{\boldsymbol{u}}^{(r)};\mu
     \bigr)
  \right\rangle_{M_u},
  \label{eq:galerkin_linear_definition}\\
  [H_r(\mu)]_{ijk}
  &=
  -\left\langle
    \boldsymbol{\phi}^{u,(r)}_i,\,
    \mathcal{N}_{h}
     \bigl(
       \boldsymbol{\phi}^{u,(r)}_j,
       \boldsymbol{\phi}^{u,(r)}_k;\mu
     \bigr)
  \right\rangle_{M_u},
  \label{eq:galerkin_quadratic_definition}\\
  [C_r^p(\mu)]_{i\ell}
  &=
  -\left\langle
    \boldsymbol{\phi}^{u,(r)}_i,\,
    \nabla_{h}(\mu)\phi^{p,(r)}_\ell
  \right\rangle_{M_u}.
  \label{eq:galerkin_pressure_coupling_definition}
\end{align}
For a third-order tensor $H_r(\mu)$, the double contraction is defined as
\begin{equation}
  \left[
    H_r(\mu):
    (\boldsymbol{a}\otimes\boldsymbol{a})
  \right]_i
  =
  \sum_{j=1}^{d_u}
  \sum_{k=1}^{d_u}
  [H_r(\mu)]_{ijk}a_j a_k.
  \label{eq:galerkin_tensor_contraction}
\end{equation}
Accordingly, the regime-local Galerkin contribution to the reduced velocity
dynamics is
\begin{equation}
  \mathscr{F}^{\mathrm{G}}_r
  (\boldsymbol{a},\boldsymbol{b};\mu)
  :=
  \boldsymbol{c}_r(\mu)
  +A_r(\mu)\boldsymbol{a}
  +H_r(\mu):
   (\boldsymbol{a}\otimes\boldsymbol{a})
  +C_r^p(\mu)\boldsymbol{b}.
  \label{eq:chart_local_galerkin_rhs_definition}
\end{equation}
The pressure-coupling term is retained because the finite-volume velocity
POD basis is not assumed to be discretely divergence-free.  All quadrature
rules, boundary contributions, and degree-of-freedom orderings entering the
reduced operators are inherited consistently from the full-order
finite-volume discretization.  When a particular full-order operator is
independent of $\mu$, the corresponding parameter argument in the reduced
operator may be omitted.

The pressure is not advanced through an independent evolution equation.
Instead, a pressure approximation is obtained from the incompressibility
constraint and the momentum balance.  For reference, in the
constant-viscosity continuous setting, for which
$\nabla\cdot\Delta\boldsymbol{u}=0$, the pressure satisfies
\begin{equation}
  -\Delta p^{\circ}
  =
  \nabla\cdot
  \left[
    (\boldsymbol{u}\cdot\nabla)\boldsymbol{u}
    -\boldsymbol{f}
  \right],
  \qquad
  \int_{\Omega}p^{\circ}\,\mathrm{d}\Omega=0,
  \label{eq:pressure_poisson_strong_form}
\end{equation}
together with the pressure-flux boundary condition induced by the momentum
equation.  At the discrete level, the Pressure--Poisson equation is assembled
using the same finite-volume gradient, divergence, quadrature, and boundary
fluxes as the full-order pressure-correction scheme.  We denote its
gauge-fixed bilinear form by
\begin{equation}
  \mathscr{A}^{PP}_h(q,s;\mu)
  :=
  \left\langle
    \nabla_h(\mu)q,\,
    \nabla_h(\mu)s
  \right\rangle_{\Omega,h},
  \label{eq:pressure_poisson_bilinear_form}
\end{equation}
where $\langle\cdot,\cdot\rangle_{\Omega,h}$ denotes the finite-volume
quadrature associated with the discrete pressure equation.

Let
$\mathscr{R}^{PP}_h(q;\boldsymbol{u},\mu)$ denote the complete discrete
Pressure--Poisson right-hand-side functional, including all boundary-flux
and lifting contributions.  Its standard volume contribution may be written
as
\begin{equation}
  \mathscr{R}^{PP}_{h,\Omega}
  (q;\boldsymbol{u},\mu)
  =
  -\left\langle
    \nabla_h(\mu)q,\,
    \mathcal{N}_{h}
     (\boldsymbol{u},\boldsymbol{u};\mu)
    -\boldsymbol{f}_{h}(\mu)
    -\mathcal{L}_{h}(\mu)\boldsymbol{u}
  \right\rangle_{\Omega,h}.
  \label{eq:pressure_poisson_weak_rhs}
\end{equation}
Terms whose discrete divergence vanishes identically may, of course, be
removed during assembly.  Retaining the functional notation
$\mathscr{R}^{PP}_h$ avoids imposing such cancellations before the boundary
conditions and discrete operators have been specified.

Testing the discrete Pressure--Poisson equation with
$q=\phi^{p,(r)}_\ell$ and inserting the local velocity and pressure
expansions from \eqref{eq:local_pod_expansions_componentwise} gives
\begin{equation}
  K_r^p(\mu)\boldsymbol{b}^{PP,(r)}
  =
  \boldsymbol{d}_r(\mu)
  +B_r(\mu)\boldsymbol{a}^{(r)}
  +Q_r(\mu):
   \bigl(
     \boldsymbol{a}^{(r)}
     \otimes
     \boldsymbol{a}^{(r)}
   \bigr).
  \label{eq:reduced_pressure_poisson_system}
\end{equation}
Here $\boldsymbol{b}^{PP,(r)}$ denotes the pressure coefficients supplied by
the reduced Pressure--Poisson model, which are distinguished from the
projected pressure coefficients $\boldsymbol{b}^{(r)}$ obtained directly
from a full-order snapshot.  The reduced pressure matrix is
\begin{equation}
  [K_r^p(\mu)]_{\ell m}
  =
  \mathscr{A}^{PP}_h
  \bigl(
    \phi^{p,(r)}_\ell,
    \phi^{p,(r)}_m;\mu
  \bigr).
  \label{eq:pressure_poisson_stiffness}
\end{equation}
To define the polynomial right-hand side, we use the exact decomposition
\begin{align}
  &\mathscr{R}^{PP}_h
  \left(
    q;\,
    \overline{\boldsymbol{u}}^{(r)}
    +\sum_{j=1}^{d_u}
     a_j\boldsymbol{\phi}^{u,(r)}_j,\,
    \mu
  \right)
  \nonumber\\
  &\qquad
  =
  \mathscr{R}^{PP}_h
  \left(
    q;\overline{\boldsymbol{u}}^{(r)},\mu
  \right)
  +\sum_{j=1}^{d_u}
   a_j
   \mathscr{R}^{PP}_{h,\mathrm{lin}}
   \left(
     q;\overline{\boldsymbol{u}}^{(r)},
     \boldsymbol{\phi}^{u,(r)}_j,\mu
   \right)
  \nonumber\\
  &\qquad\quad
  +\sum_{j=1}^{d_u}
   \sum_{k=1}^{d_u}
   a_j a_k
   \mathscr{R}^{PP}_{h,\mathrm{quad}}
   \left(
     q;
     \boldsymbol{\phi}^{u,(r)}_j,
     \boldsymbol{\phi}^{u,(r)}_k,\mu
   \right).
  \label{eq:pressure_poisson_rhs_decomposition}
\end{align}
The constant, linear, and quadratic reduced coefficients are therefore
defined by
\begin{align}
  [\boldsymbol{d}_r(\mu)]_\ell
  &=
  \mathscr{R}^{PP}_h
  \bigl(
    \phi^{p,(r)}_\ell;
    \overline{\boldsymbol{u}}^{(r)},\mu
  \bigr)
  -
  \mathscr{A}^{PP}_h
  \bigl(
    \phi^{p,(r)}_\ell,
    \overline{p}^{(r),\circ};\mu
  \bigr),
  \label{eq:pressure_poisson_constant}\\
  [B_r(\mu)]_{\ell j}
  &=
  \mathscr{R}^{PP}_{h,\mathrm{lin}}
  \bigl(
    \phi^{p,(r)}_\ell;
    \overline{\boldsymbol{u}}^{(r)},
    \boldsymbol{\phi}^{u,(r)}_j,\mu
  \bigr),
  \label{eq:pressure_poisson_linear}\\
  [Q_r(\mu)]_{\ell jk}
  &=
  \mathscr{R}^{PP}_{h,\mathrm{quad}}
  \bigl(
    \phi^{p,(r)}_\ell;
    \boldsymbol{\phi}^{u,(r)}_j,
    \boldsymbol{\phi}^{u,(r)}_k,\mu
  \bigr).
  \label{eq:pressure_poisson_quadratic}
\end{align}
The linear functional contains the contributions that are linear in the
velocity fluctuation, including the two mean--fluctuation convection terms
and any linear momentum contribution retained by the discrete
Pressure--Poisson equation.  The quadratic functional contains the
mode--mode convection contribution.  Both functionals include the boundary
terms generated by the full-order discretization.  This formulation avoids
discarding terms that vanish only under additional assumptions, such as an
exactly solenoidal velocity basis or homogeneous idealized boundary
conditions.

Because the pressure mean and all pressure POD modes belong to the
gauge-fixed space $\mathcal{Q}_{0,h}$, the constant-pressure nullspace has
already been removed.  On a connected domain, the reduced matrix
$K_r^p(\mu)$ is therefore positive definite under the standard coercivity
assumptions for the discrete Pressure--Poisson operator.  The regime-local
Pressure--Poisson map is consequently defined by
\begin{equation}
  \mathscr{Q}^{PP}_r(\boldsymbol{a};\mu)
  =
  [K_r^p(\mu)]^{-1}
  \left[
    \boldsymbol{d}_r(\mu)
    +B_r(\mu)\boldsymbol{a}
    +Q_r(\mu):
     (\boldsymbol{a}\otimes\boldsymbol{a})
  \right].
  \label{eq:chart_local_pressure_poisson_map}
\end{equation}
In implementation, \eqref{eq:chart_local_pressure_poisson_map} is evaluated
by solving the reduced linear system rather than by forming the matrix
inverse explicitly.  The map is assembled offline and remains fixed during
the training of the regime-specific learned model.  Any learned pressure
correction introduced later is defined relative to this algebraic baseline
and does not alter the underlying reduced Pressure--Poisson operator.

\subsection{Complete regime-local ROM chart}
\label{subsec:complete_regime_local_rom_chart}

The objects introduced above define the reduced representation associated
with regime $r$.  We collect them in the regime-local ROM chart
\begin{equation}
  \mathcal{C}^{\mathrm{ROM}}_r
  =
  \left(
    \mathcal{M}_r,\,
    \Gamma_p,\,
    \overline{\boldsymbol{u}}^{(r)},\,
    \Phi_u^{(r)},\,
    \overline{p}^{(r),\circ},\,
    \Phi_p^{(r)},\,
    \boldsymbol{\mu}_a^{(r)},\,
    \boldsymbol{\sigma}_a^{(r)},\,
    \boldsymbol{\mu}_b^{(r)},\,
    \boldsymbol{\sigma}_b^{(r)},\,
    \mathscr{F}^{\mathrm{G}}_r,\,
    \mathscr{Q}^{PP}_r
  \right).
  \label{eq:local_chart_contract}
\end{equation}
Thus, $\mathcal{C}^{\mathrm{ROM}}_r$ specifies the parameter region represented
by the chart, the pressure gauge, the regime-local affine velocity and
pressure spaces, the coefficient standardization, and the projected
momentum and Pressure--Poisson operators.  The chart contains only the
reduced spatial representation and the associated physics-based operators.
History-dependent features, phase variables, learned closure maps,
regime-specific time-stepping rules, and routing variables are introduced
separately when the local dynamical specialists are defined in
Section~\ref{sec:hprs_moe_specialists}.

Although the local charts may use the same reduced dimensions $d_u$ and
$d_p$, their means and POD bases are constructed from different
regime-specific training populations.  Equality of the coefficient-vector
dimensions therefore does not imply a common reduced coordinate system, and
modal coefficients from different charts cannot be identified or combined
componentwise without an explicit change of coordinates.

\subsection{Coordinate transformations between adjacent local POD spaces}
\label{subsec:adjacent_chart_interfaces}

Coordinate transformations are introduced only between dynamically adjacent
regimes.  The set of directed adjacent pairs is
\begin{equation}
  \mathcal{E}_{\mathrm{adj}}
  =
  \{
    S\!\to\!H,\,
    H\!\to\!S,\,
    H\!\to\!P,\,
    P\!\to\!H
  \}.
  \label{eq:adjacent_directed_edges}
\end{equation}
No direct transformation between the steady and periodic charts is used.

Let $(r,k)\in\mathcal{E}_{\mathrm{adj}}$.  A state represented by
$\boldsymbol{a}^{(r)}$ is first reconstructed in the physical velocity space
and then orthogonally projected onto the affine POD space of chart $k$.  The
resulting target-chart coefficient vector is
\begin{equation}
  \boldsymbol{a}^{(k\leftarrow r)}
  =
  T^{u}_{k\leftarrow r}\boldsymbol{a}^{(r)}
  +\boldsymbol{c}^{u}_{k\leftarrow r},
  \label{eq:velocity_affine_map}
\end{equation}
where
\begin{equation}
  T^{u}_{k\leftarrow r}
  =
  {\Phi_u^{(k)}}^{\mathsf T}
  M_u
  \Phi_u^{(r)},
  \qquad
  \boldsymbol{c}^{u}_{k\leftarrow r}
  =
  {\Phi_u^{(k)}}^{\mathsf T}
  M_u
  \left(
    \overline{\boldsymbol{u}}^{(r)}
    -\overline{\boldsymbol{u}}^{(k)}
  \right).
  \label{eq:velocity_map_components}
\end{equation}
The notation $\boldsymbol{a}^{(k\leftarrow r)}$ emphasizes that the mapped
coefficient vector is the best $M_u$-orthogonal representation, in chart
$k$, of the field reconstructed from chart $r$.  The transformation is exact
only when that reconstructed field belongs to the target affine POD space.

The corresponding pressure transformation is
\begin{equation}
  \boldsymbol{b}^{(k\leftarrow r)}
  =
  T^{p}_{k\leftarrow r}\boldsymbol{b}^{(r)}
  +\boldsymbol{c}^{p}_{k\leftarrow r},
  \label{eq:pressure_affine_map}
\end{equation}
with
\begin{equation}
  T^{p}_{k\leftarrow r}
  =
  {\Phi_p^{(k)}}^{\mathsf T}
  M_p
  \Phi_p^{(r)},
  \qquad
  \boldsymbol{c}^{p}_{k\leftarrow r}
  =
  {\Phi_p^{(k)}}^{\mathsf T}
  M_p
  \left(
    \overline{p}^{(r),\circ}
    -\overline{p}^{(k),\circ}
  \right).
  \label{eq:pressure_map_components}
\end{equation}
Since both local pressure means and all pressure POD modes satisfy the same
gauge condition, the mapped pressure field remains in
$\mathcal{Q}_{0,h}$.

The affine offsets in
\eqref{eq:velocity_map_components} and
\eqref{eq:pressure_map_components} account exclusively for differences
between the local means.  Consequently, fluctuations, increments, and time
derivatives are transferred by the linear parts $T^u_{k\leftarrow r}$ and
$T^p_{k\leftarrow r}$, without the affine offsets.  These transformations
are assembled offline from the local means, POD bases, and quadrature
matrices; no matrix inversion or trainable mapping is introduced.

The compatibility of two adjacent local spaces is assessed using both
basis-level and state-dependent diagnostics.  Because the two POD bases are
orthonormal with respect to the same physical inner product, the singular
values of $T^u_{k\leftarrow r}$ are the cosines of the principal angles
between the two velocity POD subspaces:
\begin{equation}
  \cos\theta_i^u(r,k)
  =
  \sigma_i
  \left(
    T^u_{k\leftarrow r}
  \right),
  \qquad
  i=1,\ldots,d_u.
  \label{eq:velocity_principal_angles}
\end{equation}
An analogous relation holds for the pressure transformation.  Singular
values close to one indicate directions that are well represented in both
subspaces, whereas small singular values identify directions that are
poorly shared across the overlap.

For a velocity field $\boldsymbol{u}$, the irreducible relative projection
error associated with representation in chart $k$ is
\begin{equation}
  e_{\mathrm{proj}}^u(\boldsymbol{u};k)
  =
  \frac{
    \left\|
      \left(
        I
        -\Phi_u^{(k)}
         {\Phi_u^{(k)}}^{\mathsf T}
         M_u
      \right)
      \left(
        \boldsymbol{u}
        -\overline{\boldsymbol{u}}^{(k)}
      \right)
    \right\|_{M_u}
  }{
    \left\|
      \boldsymbol{u}
      -\overline{\boldsymbol{u}}^{(k)}
    \right\|_{M_u}
    +\varepsilon
  }.
  \label{eq:target_projection_floor}
\end{equation}
This quantity is a representation error of the target affine POD space and
therefore provides a lower bound for any prediction restricted to that
space.  The corresponding pressure projection error is defined analogously
after applying the gauge map $\Gamma_p$.

Because the transformations are projections between generally distinct
subspaces, their composition is not expected to be the identity.  For
example, the velocity coefficients obtained after a round trip
$r\to k\to r$ are
\begin{equation}
  \boldsymbol{a}^{(r\leftarrow k\leftarrow r)}
  =
  T^u_{r\leftarrow k}
  \left(
    T^u_{k\leftarrow r}\boldsymbol{a}^{(r)}
    +\boldsymbol{c}^u_{k\leftarrow r}
  \right)
  +\boldsymbol{c}^u_{r\leftarrow k}.
  \label{eq:velocity_round_trip_map}
\end{equation}
A field-level round-trip error may then be defined by
\begin{equation}
  e_{\mathrm{rt}}^u
  \left(
    \boldsymbol{a}^{(r)};r,k
  \right)
  =
  \frac{
    \left\|
      \mathcal{D}_r^u
      \left(
        \boldsymbol{a}^{(r\leftarrow k\leftarrow r)}
      \right)
      -
      \mathcal{D}_r^u
      \left(
        \boldsymbol{a}^{(r)}
      \right)
    \right\|_{M_u}
  }{
    \left\|
      \mathcal{D}_r^u
      \left(
        \boldsymbol{a}^{(r)}
      \right)
      -
      \overline{\boldsymbol{u}}^{(r)}
    \right\|_{M_u}
    +\varepsilon
  }.
  \label{eq:velocity_round_trip_error}
\end{equation}
Projection errors, round-trip errors, principal angles, and the singular
spectra of the cross-Gram matrices are used to assess representational
compatibility between adjacent charts.  They are diagnostic quantities and
are not included as trainable penalties once the local POD spaces and
specialists have been fixed.

When a regime-specific dynamical model uses a sequence of past states, the
coordinate transformation is applied independently to each unstandardized
POD coefficient vector, while the associated physical timestamps are
retained unchanged.  The mapped sequence is then standardized using the
statistics of the target chart, and all target-specific temporal features,
including finite differences, phase variables, and closure histories, are
reconstructed according to the definition of the target specialist.  A
standardized feature vector from the source chart is therefore never
transferred directly.

The adjacent transformations provide consistent initialization and
compatibility diagnostics for the local models; they do not define a global
POD space or a common reduced vector field.  In particular, modal
coefficients from distinct charts are not combined componentwise.  Any
admissible soft combination introduced in
Section~\ref{sec:hierarchical_ral_moe} is performed only after the local
predictions have been reconstructed in the common physical velocity and
pressure spaces.


